\section{Two coordinate systems, two sign rules}
\label{sec:ssrsigns}

Before any rule can be written down, one pair of conventions must be
fixed, because it is where nearly every practical error in smile
transport lives.  The market trades \emph{strikes}: dollar levels,
printed on contracts, unchanged when the underlying moves.  The book
draws \emph{smiles}: curves in log-moneyness $k=\log(K/F)$, measured
against the forward.  When the forward moves, every strike stays put
and every label changes.  Two coordinate systems, one of which just
slid under the other.

Fix the conventions once, and hold them for the whole chapter.

\begin{definition}[Sign conventions]
\label{def:ssrsigns}
(i)~The move is $H=\log(F_{\mathrm{new}}/F_{\mathrm{old}})$: positive
when the forward rises.  (ii)~Log-moneyness $k=\log(K/F)$ is always
measured against the \emph{prevailing} forward: today's curve uses
today's forward, tomorrow's curve uses tomorrow's.  (iii)~A
transported \emph{curve} is a function of new moneyness; a transported
\emph{quote} --- a fixed strike --- is a point whose moneyness label
changes.
\end{definition}

One wording convention travels with (i), stated once for the whole
chapter: a ``four percent move'' always means the \emph{log} move
$H=\pm0.04$ --- so a $4\%$ rise lands the forward at
$100\,e^{0.04}=104.08$, not $104$ --- and, likewise, a strike ``ten
percent below the money'' means $k=-0.10$.

Now do the two computations these conventions force, by hand.  Let
today's forward be $F=100$ and let it rise
$\MacSsrSignMovePct\%$, so $H=+0.04$ and
$F_{\mathrm{new}}=100\,e^{0.04}=104.08$.

First, follow one \emph{quote}.  The strike $K=100$ has today's label
$k=\log(100/100)=0$: it is at the money.  Tomorrow its label is
\[
k_{\mathrm{new}}
=\log\!\frac{100}{104.08}
=\log\!\frac{K}{F_{\mathrm{old}}e^{H}}
=k-H=-0.04 .
\]
A fixed strike's label slides \emph{down} by the move: the
\textbf{quote rule}, $k\mapsto k-H$.  This is just relabeling; no
price and no volatility has been touched.

Second, build the new \emph{curve}, under the simplest hypothesis a
desk uses: every fixed strike keeps yesterday's volatility
(``sticky-strike'').  To evaluate the new curve at new label $k$, ask
which strike now sits there.  It is $K=F_{\mathrm{new}}e^{k}$, and its
\emph{old} label was
\[
\log\!\frac{K}{F_{\mathrm{old}}}
=\log\!\frac{F_{\mathrm{new}}e^{k}}{F_{\mathrm{old}}}
=k+H .
\]
That strike's volatility is unchanged, so the new curve reads the old
curve there:
\begin{equation}
\sigma_{\mathrm{new}}(k)=\sigma_{\mathrm{old}}(k+H)
\qquad\text{(sticky-strike)} .
\label{eq:ssrcurverule}
\end{equation}
The \textbf{curve rule} reads \emph{forward}, at $k+H$ --- a plus ---
while the quote rule slides labels to $k-H$ --- a minus.  Both signs
are correct, and they point in opposite directions in $k$.  This is
the trap.  Swap either sign and the transported smile's skew response
flips; the error is silent, because the resulting picture still looks
like a plausible smile.  The rule is to derive each sign the way
we just did --- ``which strike sits here now?''\ for curves, ``where
did my strike's label go?''\ for quotes --- and never to guess.

\begin{figure}[htbp]
\centering
\paperfig{\textwidth}{fig_ssr_signs.pdf}
\caption{The two signs on one $+\MacSsrSignMovePct\%$ move (frozen
hero smile, sticky-strike).  The fixed strike whose fitted volatility
is \MacSsrSignQuoteVolPct\% keeps that volatility while its label
slides
left to $k-H$ (filled to open circle: the quote rule).  The new curve
(blue) at any label $k$ reads the old curve (orange) at $k+H$ (filled
to open square: the curve rule).  Opposite directions, both correct.}
\label{fig:ssrsigns}
\end{figure}

\Cref{fig:ssrsigns} draws both rules on the frozen node's smile under
the same $+\MacSsrSignMovePct\%$ move.  The orange curve is today's
smile; the blue curve is the sticky-strike tomorrow, the whole shape
carried left by the re-indexing \eqref{eq:ssrcurverule}.  Follow the
circles first: the strike carrying \MacSsrSignQuoteVolPct\% sits at
label $+0.02$ today (filled circle) and at $-0.02$ tomorrow (open
circle), at exactly the same height --- same volatility, new label,
the quote rule.  Now follow the squares: to know tomorrow's curve at
label $-0.08$, walk \emph{right} by $H$ to the old curve at $-0.04$
(filled square) and carry that value across (open square) --- the
curve rule.  The circles move one way, the squares' reading goes the
other, and the figure is worth returning to whenever a sign feels
doubtful.

One more reading of \eqref{eq:ssrcurverule} is worth a sentence,
because it explains the name of the whole subject.  Under
sticky-strike the curve slides along the strike axis; under the
opposite pure hypothesis --- the smile is a fixed function of
moneyness, ``sticky-moneyness'' --- the curve would not move at all in
$k$.  Real desks argue precisely about how sticky the smile is in
which coordinate.  The conventions are now fixed; the next section
shows that, once shapes are preserved, the entire argument collapses
to a single number.
